% Arquivo principal: e o unico que se compila. O texto fica em sections/,
% um arquivo por secao, incluido com \input na ordem do artigo.
\documentclass[conference]{IEEEtran}
% \IEEEoverridecommandlockouts % so necessario se usar \thanks no titulo (financiamento)
\usepackage[T1]{fontenc}
\usepackage[english,brazilian]{babel} % brazilian (o ultimo) e o idioma principal
\usepackage{cite}
\usepackage{amsmath,amssymb,amsfonts}
\usepackage{algorithmic}
\usepackage{graphicx}
\graphicspath{{../}} % figuras ficam em ../images/ -> \includegraphics{images/fig.png}
\usepackage{textcomp}
\usepackage{xcolor}
\usepackage{url}
\usepackage{array}
\usepackage{tikz}
\usetikzlibrary{positioning,arrows.meta,fit}
\usepackage{listings}
\lstset{basicstyle=\scriptsize\ttfamily, breaklines=true, columns=fullflexible, keepspaces=true,
  frame=single, framesep=2pt, aboveskip=2pt, belowskip=6pt, breakindent=0pt}

\renewcommand{\IEEEkeywordsname}{Palavras-chave} % o babel nao traduz o "Index Terms" do IEEEtran

\begin{document}

\title{Contratos de Handoff, Propagação e Atribuição de Erros em Pipelines de IA Generativa
% \thanks{Agência de fomento, se houver.}
}

\author{
\IEEEauthorblockN{Victor Amarante F. da Silva}
\IEEEauthorblockA{\textit{Centro de Informática} \\
\textit{Universidade Federal de Pernambuco}\\
Recife, Brasil \\
vafs@cin.ufpe.br}
\and
\IEEEauthorblockN{Cleber Zanchettin}
\IEEEauthorblockA{\textit{Orientador} \\
\textit{Centro de Informática} \\
\textit{Universidade Federal de Pernambuco}\\
Recife, Brasil \\
cz@cin.ufpe.br}
}

\maketitle

% !TEX root = ../tcc.tex
% Resumo: ate ~250 palavras, sem citacoes nem notas de rodape (regra IEEE).

\begin{abstract}
Em pipelines de IA generativa com vários estágios, a saída de um estágio vira a entrada do próximo, e um erro introduzido cedo pode chegar à resposta final sem que se saiba onde começou.
Este trabalho investiga como o contrato de handoff, o formato da informação trocada entre estágios, afeta a propagação, a correção e a atribuição de origem desse erro.
Um pipeline fixo de três estágios (Extrator, Analista e Redator) roda sobre o FinanceBench com quatro contratos: texto livre (A), tipado (B), tipado com lineage (C) e tipado com lineage e o trecho de evidência da fonte (D).
Um injetor determinístico corrompe um claim por execução, e o mesmo estado corrompido é projetado nas quatro condições.
O protocolo foi repetido com três modelos: GPT-5.6 Luna, Claude Haiku 4.5 e o SLM Qwen3.5 9B.
Nos três, a estrutura sozinha não melhorou a acurácia final e aumentou a profundidade de propagação.
O lineage elevou a atribuição correta do claim de origem de menos de 0,09 para cerca de 0,4, mas não reduziu a sobrevivência do erro.
Nos dois modelos de API, só a evidência reduziu a sobrevivência e a propagação do erro, e a atribuição chegou a cerca de 0,8 (+0,38 a +0,43 sobre C), com quase quatro vezes mais tokens que o texto livre.
No Qwen, a evidência melhorou a acurácia final, mas não a contenção nem a atribuição.
Estrutura e lineage, sozinhos, não tornam o pipeline mais confiável: o ganho vem da evidência que permite conferir a proveniência, e depende de o modelo saber usá-la.
\end{abstract}

\begin{IEEEkeywords}
pipelines de IA generativa, propagação de erros, sistemas multiagente, observabilidade, atribuição de causa raiz, contratos de handoff
\end{IEEEkeywords}

% !TEX root = ../tcc.tex
\section{Introdução}\label{sec:introducao}

% - Contexto: pipelines multiestágio de LLMs e handoffs entre componentes probabilísticos.
% - Problema: um erro introduzido cedo pode ser preservado, transformado, amplificado,
%   corrigido ou contido -- e hoje é difícil observar e atribuir sua origem.
% - Foco: observabilidade e atribuição no nível do sistema/harness, não do modelo.
% - Pergunta central e as perguntas de pesquisa:
%     RQ1 (A x B): efeito da estruturação do handoff;
%     RQ2 (B x C): efeito do lineage/proveniência;
%     RQ3 (C x D): efeito adicional da evidência da fonte.
% - Contribuições (lista com \begin{itemize}).
% - Organização do artigo (Seção~\ref{sec:trabalhos} ... Seção~\ref{sec:conclusao}).

% !TEX root = ../tcc.tex
\section{Trabalhos Relacionados}\label{sec:trabalhos}

% Sugestão de agrupamento (chaves em referencias.bib):
% - Propagação de erros em pipelines multiagente: \cite{singh2026snowball, barrak2025traceability}
% - Saídas estruturadas e validade x correção: \cite{geng2025structured, ray2026constraint}
% - Harness, observabilidade e diagnóstico de falhas: \cite{zhong2026harness, chen2026failed, ran2026caveagent}
% - Arquiteturas de agentes generativos (contexto): \cite{park2023generative, altera2024sid, park2024selfreports}
% Fechar com a lacuna que este trabalho preenche.

% !TEX root = ../tcc.tex
\section{Metodologia}\label{sec:metodologia}

\subsection{Pipeline}\label{sec:pipeline}
% Documento-fonte -> Extrator -> Analista -> Redator -> Resposta final.
% Modelo, prompts, documentos e lógica dos estágios constantes entre condições.

\subsection{Unidade de análise: o claim}\label{sec:claim}
% claim_id, derived_from e a cadeia de proveniência (ex.: C31 -> C17 -> C08).

\subsection{Condições de contrato de handoff}\label{sec:condicoes}
% A (texto livre), B (tipado), C (tipado + lineage), D (tipado + lineage + evidência).
% Tabela comparativa e o controle de validade interna: em C, source_id/page são
% só metadados; apenas D recebe o trecho de evidência.

\subsection{Injeção de erros}\label{sec:injecao}
% Três classes (valor, percentual, metadado), perturbações determinísticas
% registradas como ground truth; mesmo estado corrompido serializado em A/B/C/D.

\subsection{Classificação do destino do erro}\label{sec:destino}
% Preservado / transformado / amplificado / corrigido / contido.

% !TEX root = ../tcc.tex
\section{Configuração Experimental}\label{sec:configuracao}

\subsection{Conjunto de dados}\label{sec:dados}
% FinanceBench \cite{islam2023financebench}: subconjunto usado, tipos de pergunta.

\subsection{Modelos}\label{sec:modelos}
% Famílias de modelo avaliadas e configuração (temperatura, seeds, execução local/API).

\subsection{Métricas}\label{sec:metricas}
% Correção: Final Factual Accuracy, Correction Rate.
% Contenção: Error Survival Rate (ESR), Propagation Depth, Error Amplification.
% Observabilidade de causa raiz: RCCA, FESA.
% Custo operacional: tokens, latência, chamadas, custo monetário.
% Schema Validity (B/C/D), separada de correção semântica.

% !TEX root = ../tcc.tex
\section{Resultados}\label{sec:resultados}

\subsection{RQ1: efeito da estruturação (A $\times$ B)}\label{sec:rq1}

\subsection{RQ2: efeito do lineage (B $\times$ C)}\label{sec:rq2}

\subsection{RQ3: efeito da evidência (C $\times$ D)}\label{sec:rq3}

\subsection{Custo operacional}\label{sec:custo}

\subsection{Replicação entre famílias de modelo}\label{sec:replicacao}

% Padrão IEEE: legenda da TABELA acima, legenda da FIGURA abaixo; \label
% sempre depois do \caption. Citar no texto como Tabela~\ref{tab:...} e Fig.~\ref{fig:...}.
%
% \begin{table}[htbp]
% \caption{Resultados por condição}
% \label{tab:resultados}
% \centering
% \begin{tabular}{|l|c|c|c|c|}
% \hline
% \textbf{Métrica} & \textbf{A} & \textbf{B} & \textbf{C} & \textbf{D} \\
% \hline
% ESR & & & & \\
% RCCA & & & & \\
% \hline
% \end{tabular}
% \end{table}
%
% \begin{figure}[htbp]
% \centering
% \includegraphics[width=\columnwidth]{images/fig1.png}
% \caption{Legenda da figura.}
% \label{fig:exemplo}
% \end{figure}

% !TEX root = ../tcc.tex
\section{Discussão}\label{sec:discussao}

% Interpretação dos resultados por mecanismo (estrutura, lineage, evidência),
% resultados negativos e implicações práticas para quem projeta pipelines.

% !TEX root = ../tcc.tex
\section{Ameaças à Validade}\label{sec:ameacas}

As ameaças estão organizadas segundo a classificação usual em experimentação de software, que separa validade interna, de construto, externa e de conclusão \cite{wohlin2012experimentation}.

\subsection{Validade interna}\label{sec:ameacas-interna}

A ameaça mais séria à comparação entre condições é o vazamento de informação pelo identificador da fonte.
Na condição C, o \texttt{source\_id} é o nome do documento no FinanceBench, e em 128 das 142 perguntas do braço de metadado no GPT esse nome contém o ano fiscal original, que contradiz o ano corrompido.
A condição C deveria carregar apenas proveniência, mas nesse braço carregou também um pedaço de evidência.
O efeito da RQ2 no erro de metadado está, por isso, inflado, e a estimativa mais limpa do efeito do lineage é a do braço de valor, em que o identificador não revela nada sobre o valor corrompido.
A correção natural é repetir a condição C com identificadores opacos.

Os modelos não rodaram com a mesma configuração de raciocínio. O GPT-5.6 Luna usou esforço \emph{medium}, enquanto Haiku e Qwen rodaram sem raciocínio.
Como a configuração é constante dentro de cada modelo, isso não afeta as comparações entre condições, que são o objeto das RQs, mas limita as comparações entre modelos.
Pelo mesmo motivo, cada execução rodou com uma única seed e com a temperatura padrão de cada provedor.
A variação entre execuções repetidas foi medida no piloto apenas para o GPT, em que duas seeds deram respostas idênticas, e não foi medida para Haiku e Qwen.
Essa variação aleatória afeta igualmente as quatro condições, então tende a alargar os intervalos, e não a criar diferenças entre elas.

Por fim, as condições não diferem só na informação que carregam, mas também no tamanho das mensagens, já que um handoff de D tem quase quatro vezes mais tokens que um de A.
Essa diferença é inerente ao que se quer medir e não pode ser separada da informação adicional sem outro desenho, como condições de controle com texto de enchimento.

\subsection{Validade de construto}\label{sec:ameacas-construto}

A RCCA aceita dois caminhos de acerto, e só as condições C e D têm os dois.
Em A e B, o auditor só acerta se informar a métrica e o valor corrompido. Em C e D, ele também acerta ao apontar o identificador do claim, e na condição C cerca de 0,3 dos aproximadamente 0,4 de RCCA vêm desse caminho.
A Seção~\ref{sec:rq2} mostra a RCCA recalculada com o mesmo critério nas quatro condições, e nela o ganho de B para C quase desaparece.
As duas medidas respondem a perguntas diferentes, a de qual claim originou o erro e a de qual valor está errado, e os resultados devem ser lidos com essa distinção em mente.

O auditor usa o mesmo modelo do pipeline, e a sua confiabilidade varia com o modelo.
No Haiku e no Qwen, ele apontou culpados em 30\% a 74\% das execuções sem erro injetado, o que torna a FESA desses modelos ininterpretável e limita também a RCCA.
Um auditor fixo, mais forte e igual para os três experimentos separaria melhor a capacidade de atribuição do contrato da capacidade do modelo que audita.

O instrumento de medida trabalha com os números presentes no texto de cada estágio.
Erros expressos por extenso ou números que coincidem por acaso com o valor procurado podem escapar ou ser contados indevidamente, e o ano original pode reaparecer na resposta por estar na própria pergunta, o que motivou a regra mais estrita para distinguir correção de contenção (Seção~\ref{sec:destino}).
A acurácia final também depende de casar números com a resposta do benchmark, e as 18 a 22 perguntas sem resposta numérica ficaram fora dessa métrica.

Os erros injetados são sintéticos.
A transposição de dígitos e o deslocamento de um ano produzem erros plausíveis, mas não necessariamente parecidos com as alucinações que um modelo comete sozinho, e Singh e Pawar usaram perturbações de outro tipo, com deslocamentos multiplicativos de 15\% a 40\% \cite{singh2026snowball}.
A facilidade de detectar o erro de ano, que pode entrar em conflito com a pergunta, também difere da do erro de valor, o que justifica reportar os dois braços separadamente.

\subsection{Validade externa}\label{sec:ameacas-externa}

O experimento cobre um único domínio, finanças, com um único benchmark, a amostra aberta de 150 perguntas do FinanceBench \cite{islam2023financebench}.
O documento-fonte é o trecho de evidência já recortado pelo benchmark, e não o documento completo, então o experimento não inclui recuperação nem documentos longos, situações em que a evidência carregada pelo contrato pode ser mais difícil de localizar e de usar.
A topologia é linear e tem três estágios, enquanto sistemas reais podem envolver conversas entre vários agentes com ferramentas \cite{wu2023autogen} ou simulações com mais de mil agentes \cite{altera2024sid}.
Os três modelos são das faixas mais baratas de cada família, e não se sabe se modelos maiores reproduziriam o mesmo padrão.

\subsection{Validade de conclusão}\label{sec:ameacas-conclusao}

A Tabela~\ref{tab:deltas} reúne 63 comparações, sem correção para comparações múltiplas.
Com intervalos de 95\%, é esperado que algumas diferenças apareçam como significativas por acaso.
Os achados centrais não dependem desse limite, porque seus intervalos ficam longe de zero e se repetem entre modelos, como o aumento da profundidade de A para B, o aumento da RCCA de B para C e os efeitos da evidência nos modelos de API.
Resultados isolados, com intervalos que mal excluem zero, devem ser lidos com cautela.
O subestudo do Qwen tem apenas 19 perguntas, e seus resultados são indicativos.

% !TEX root = ../tcc.tex
\section{Conclusão}\label{sec:conclusao}

% Resposta direta a cada RQ, contribuição principal e trabalhos futuros.

\appendices % o \appendices fica aqui porque abre os dois apendices
% !TEX root = ../tcc.tex
% Tabelas de apoio, citadas na Secao de Resultados. Saem dos mesmos scripts
% das demais (build_tables.py e build_complementares.py).
\section{Tabelas complementares}\label{sec:apendice-tabelas}

As tabelas abaixo detalham o subestudo de atribuição de estágio, e o fatorial entre lineage e evidência no Qwen3.5 9B, discutidos na Seção~\ref{sec:resultados}.

\begin{table*}[t]
\caption{Resultados do subestudo de estágio por modelo e condição (injeção no Analista)}
\label{tab:subestudo}
\centering
\footnotesize
\setlength{\tabcolsep}{3.2pt}
\begin{tabular}{|l|cccc|cccc|cccc|}
\hline
 & \multicolumn{4}{c|}{\textbf{GPT-5.6 Luna}} & \multicolumn{4}{c|}{\textbf{Claude Haiku 4.5}} & \multicolumn{4}{c|}{\textbf{Qwen3.5 9B}} \\
\textbf{Métrica} & A & B & C & D & A & B & C & D & A & B & C & D \\
\hline
\multicolumn{13}{|l|}{\textit{Correção}} \\
Acurácia factual final & 0,21 & 0,23 & 0,21 & 0,32 & 0,30 & 0,30 & 0,28 & 0,39 & 0,36 & 0,21 & 0,14 & 0,14 \\
Acurácia estrita$^{b}$ & 0,11 & 0,09 & 0,09 & 0,23 & 0,07 & 0,11 & 0,13 & 0,22 & 0,00 & 0,00 & 0,00 & 0,07 \\
Taxa de correção$^{c}$ & 0,20 & 0,26 & 0,25 & 0,27 & 0,48 & 0,52 & 0,38 & 0,45 & 0,56 & 0,29 & 0,14 & 0,29 \\
Retenção$^{d}$ & 0,34 & 0,31 & 0,31 & 0,41 & 0,50 & 0,55 & 0,55 & 0,55 & 0,50 & 0,33 & 0,17 & 0,33 \\
\multicolumn{13}{|l|}{\textit{Contenção}} \\
ESR & 0,89 & 0,89 & 0,93 & 0,71 & 0,94 & 0,87 & 0,81 & 0,83 & 0,84 & 0,84 & 0,89 & 0,63 \\
Profundidade de propagação & 0,91 & 0,91 & 0,95 & 0,82 & 0,94 & 0,91 & 0,87 & 0,81 & 0,89 & 0,95 & 0,95 & 0,89 \\
\quad só repetição literal$^{e}$ & 0,82 & 0,80 & 0,80 & 0,70 & 0,61 & 0,72 & 0,67 & 0,54 & 0,53 & 0,68 & 0,58 & 0,63 \\
\quad com valores novos$^{e}$ & 0,09 & 0,11 & 0,14 & 0,12 & 0,33 & 0,19 & 0,20 & 0,28 & 0,37 & 0,26 & 0,37 & 0,26 \\
Amplificação & 0,11 & 0,12 & 0,21 & 0,16 & 0,56 & 0,33 & 0,22 & 0,37 & 0,37 & 0,42 & 0,58 & 0,42 \\
\quad por valor emitido & 0,02 & 0,02 & 0,02 & 0,02 & 0,07 & 0,03 & 0,03 & 0,04 & 0,07 & 0,05 & 0,08 & 0,05 \\
\multicolumn{13}{|l|}{\textit{Observabilidade}} \\
RCCA$^{f}$ & 0,96 & 0,98 & 0,98 & 0,98 & 0,57 & 0,65 & 0,65 & 0,63 & 0,32 & 0,53 & 0,47 & 0,58 \\
\quad placebo$^{a}$ & 0,05 & 0,07 & 0,04 & 0,02 & 0,31 & 0,37 & 0,43 & 0,41 & 0,37 & 0,32 & 0,37 & 0,42 \\
RCCA, critério original$^{g}$ & 0,96 & 0,98 & 0,98 & 0,98 & 0,57 & 0,65 & 0,65 & 0,63 & 0,32 & 0,53 & 0,47 & 0,58 \\
FESA & 0,96 & 0,98 & 0,98 & 0,98 & 0,70 & 0,76 & 0,81 & 0,80 & 0,26 & 0,53 & 0,53 & 0,58 \\
\quad placebo$^{a}$ & 0,11 & 0,09 & 0,09 & 0,07 & 0,43 & 0,44 & 0,54 & 0,52 & 0,11 & 0,32 & 0,16 & 0,21 \\
Falso positivo do auditor$^{a}$ & 0,21 & 0,19 & 0,21 & 0,15 & 0,58 & 0,65 & 0,67 & 0,61 & 0,84 & 0,74 & 0,69 & 0,77 \\
Concordância de lineage & -- & -- & 0,89 & 0,86 & -- & -- & 0,89 & 0,89 & -- & -- & 0,95 & 0,89 \\
\multicolumn{13}{|l|}{\textit{Validade e linha de base}} \\
Validade de contrato & -- & 1,00 & 1,00 & 1,00 & -- & 1,00 & 1,00 & 1,00 & -- & 0,68 & 0,84 & 0,95 \\
Acurácia do controle$^{a}$ & 0,70 & 0,71 & 0,69 & 0,72 & 0,66 & 0,65 & 0,62 & 0,70 & 0,52 & 0,48 & 0,45 & 0,50 \\
\multicolumn{13}{|l|}{\textit{Custo (Analista + Redator)}} \\
Tokens por execução & 273 & 605 & 1.034 & 1.199 & 275 & 856 & 1.436 & 1.574 & 271 & 411 & 748 & 883 \\
\quad de saída & 71 & 168 & 266 & 321 & 63 & 100 & 149 & 182 & 48 & 82 & 148 & 171 \\
Latência (s) & 1,8 & 2,7 & 3,5 & 4,0 & 1,3 & 1,7 & 2,0 & 2,2 & 0,7 & 1,2 & 1,9 & 2,2 \\
US\$ por mil execuções & 0,13 & 0,29 & 0,47 & 0,56 & 0,53 & 1,25 & 2,03 & 2,30 & -- & -- & -- & -- \\
\hline
\end{tabular}
\\[2pt]{\scriptsize\raggedright $^{a}$Calculados sobre as execuções de controle (sem injeção); os placebos medem com que frequência o auditor do controle aponta o claim (ou o estágio) que o injetor corromperia. Demais métricas: média por pergunta sobre as execuções injetadas. $^{b}$Contém o número do benchmark e não contém o valor corrompido. $^{c}$Entre as execuções com controle correto e erro com efeito; o conjunto elegível muda com a condição (7 a 28 perguntas). $^{d}$Fração de acertos nas perguntas cujo controle acerta nas quatro condições (78, 74, 53 perguntas), sem o corte de erro com efeito. $^{e}$Estágios alterados apenas porque repetem o valor corrompido, e estágios com números novos em relação ao controle. $^{f}$Critério de conteúdo, comum às quatro condições: o auditor informa a métrica e o valor do claim-alvo como aparecem no trace. $^{g}$Critério original: identificador do claim (só em C e D) ou métrica e valor corrompido.\par}
\end{table*}

\begin{table}[t]
\caption{Fatorial entre lineage e evidência e efeito da instrução de conferência no Qwen3.5 9B (estudo principal)}
\label{tab:fatorial-qwen}
\centering
\footnotesize
\setlength{\tabcolsep}{3.2pt}
\begin{tabular}{|l|ccccc|}
\hline
 & \multicolumn{5}{c|}{\textbf{Qwen3.5 9B}} \\
\textbf{Métrica} & B & C & E & D & D$'$ \\
\multicolumn{1}{|l|}{} & \scriptsize --/-- & \scriptsize L/-- & \scriptsize --/E & \scriptsize L/E & \scriptsize L/E+i \\
\hline
Acurácia factual final & 0,30 & 0,28 & 0,36 & 0,37 & 0,37 \\
Retenção & 0,58 & 0,58 & 0,70 & 0,71 & 0,73 \\
ESR & 0,69 & 0,71 & 0,63 & 0,66 & 0,63 \\
Sobrevivência literal$^{a}$ & 0,23 & 0,25 & 0,21 & 0,11 & 0,13 \\
RCCA & 0,36 & 0,45 & 0,47 & 0,45 & 0,48 \\
FESA líquida & 0,25 & 0,23 & 0,30 & 0,33 & 0,33 \\
Tokens por execução & 946 & 1.912 & 1.334 & 2.171 & 2.232 \\
\hline
\end{tabular}
\\[4pt]
\resizebox{\columnwidth}{!}{{\scriptsize\setlength{\tabcolsep}{2pt}
\begin{tabular}{|l|ccc|}
\hline
 & \multicolumn{3}{c|}{\textbf{Qwen3.5 9B}} \\
\textbf{Diferença pareada} & B $\to$ E & E $\to$ D & D $\to$ D$'$ \\
\hline
Acurácia factual final & +0,05 [$-$0,01; +0,12] & +0,02 [$-$0,04; +0,07] & +0,00 [$-$0,05; +0,04] \\
Retenção & +0,12 [+0,00; +0,24] & +0,02 [$-$0,08; +0,11] & +0,02 [$-$0,09; +0,12] \\
ESR & $-$0,07 [$-$0,14; +0,00] & +0,03 [$-$0,03; +0,09] & $-$0,03 [$-$0,10; +0,03] \\
Sobrevivência literal$^{a}$ & $-$0,03 [$-$0,10; +0,05] & \textbf{$-$0,10 [$-$0,17; $-$0,03]} & +0,03 [$-$0,03; +0,09] \\
RCCA & \textbf{+0,11 [+0,05; +0,18]} & $-$0,02 [$-$0,10; +0,05] & +0,03 [$-$0,04; +0,09] \\
FESA líquida & +0,06 [$-$0,04; +0,16] & +0,03 [$-$0,07; +0,12] & +0,01 [$-$0,09; +0,11] \\
Tokens por execução & +388 & +837 & +61 \\
\hline
\end{tabular}}}
\\[2pt]{\scriptsize\raggedright L: lineage; E: evidência; i: instrução de conferência. Médias e diferenças por pergunta sobre as mesmas 137 perguntas, com IC95\% bootstrap; em negrito, intervalo que não inclui zero. A retenção usa as perguntas cujo controle acerta nas seis condições (A a D, E e D$'$). $^{a}$Fração das execuções do erro de valor em que o valor corrompido aparece literalmente na resposta final; não depende da comparação com o controle.\par}
\end{table}


% !TEX root = ../tcc.tex
% Gerado por scripts/build_apendice.py a partir de handoff/prompts.py. Nao edite a mao.
\section{Prompts e schemas}\label{sec:apendice}

Os prompts abaixo são os que rodaram no experimento, em inglês. A instrução de cada estágio é a de tarefa seguida da de formato da condição. A mensagem do usuário traz a pergunta e o handoff do estágio anterior, sob os rótulos \texttt{QUESTION} e \texttt{INPUT FROM <ESTÁGIO>}, e a do auditor traz a pergunta, o trace e a resposta final.

\noindent\textit{Extrator}
\begin{lstlisting}
You are the EXTRACTOR stage of a financial analysis pipeline.

Read the source excerpt and extract every atomic factual claim needed to answer the question.

Rules:
- Extract ONLY figures that appear literally in the source. Copy them exactly.
- NEVER compute, derive or aggregate. No ratios, margins, growth rates, averages,
  differences or percentages that are not printed in the source. Deriving figures is
  the ANALYST's job, not yours: extract the raw inputs it will need.
- One claim per atomic fact.
- NEVER rescale a figure. Report it in the scale the statement itself uses and pick the
  matching unit: a statement headed "(In thousands)" gives USD_K, "(In millions)" gives
  USD_M. A line reading "4,713,500" under "(In thousands)" is `value=4713500, unit=USD_K`
  -- never `value=4713.5, unit=USD_M`. Rescaling is arithmetic, and it is where unit
  errors come from.
- `unit` must be one of: USD_K, USD_M, USD_B, PCT, RATIO.
- `evidence` must be a literal substring of the source excerpt, and it must contain
  the digits of `value`.
- Extract only what the source states. If a value is absent, do not invent it.

The question tells you WHICH figures matter, never what to compute. If the question
asks for an operating margin and the statement prints "Operating income 1,306" and
"Net sales 11,503", extract those two figures as they are. Do NOT emit a claim like
`operating income margin = 11.4 PCT`: that figure is not in the source, and the
analyst computes it.
\end{lstlisting}

\noindent\textit{Analista: instrução de tarefa}
\begin{lstlisting}
You are the ANALYST stage of a financial analysis pipeline.

You receive the output of the EXTRACTOR. Reason over it to produce what is needed to answer the question: comparisons, growth rates, margins or the direct figure.

Rules:
- Work only from the input you received. You have no access to the source document.
- Do not invent facts that are not derivable from the input.
- Show the resulting figures, not the arithmetic.
\end{lstlisting}

\noindent\textit{Analista: formato da condição A}
\begin{lstlisting}
Write your analysis as plain prose.
\end{lstlisting}

\noindent\textit{Analista: formato da condição B}
\begin{lstlisting}
Return typed claims. Each claim has metric, value, unit and year.
\end{lstlisting}

\noindent\textit{Analista: formato da condição C}
\begin{lstlisting}
Return typed claims. Assign each claim you produce an id of the form A01, A02, ... and set `generated_by` to 'analyst'. In `derived_from`, list the ids of the input claims that this claim came from. Carry over `source_id` and `page` from the input claim you used. `source_id` and `page` are metadata only: you cannot look the document up.
\end{lstlisting}

\noindent\textit{Analista: formato da condição D}
\begin{lstlisting}
Return typed claims. Assign each claim you produce an id of the form A01, A02, ... and set `generated_by` to 'analyst'. In `derived_from`, list the ids of the input claims that this claim came from. Carry over `source_id`, `page` and the `evidence` snippet of the input claim you used.
\end{lstlisting}

\noindent\textit{Redator: instrução de tarefa}
\begin{lstlisting}
You are the WRITER stage of a financial analysis pipeline.

You receive the output of the ANALYST. Write the final answer for an end user.

Rules:
- Base the answer only on the input you received. You have no access to the source document.
- Be concise: one to three sentences.
- State the figure and its unit explicitly.
\end{lstlisting}

\noindent\textit{Redator: formato da condição A}
\begin{lstlisting}
Write the final answer as plain prose.
\end{lstlisting}

\noindent\textit{Redator: formato da condição B}
\begin{lstlisting}
Return `answer_text` with the final answer, plus `answer_claim`: the typed claim that the answer asserts.
\end{lstlisting}

\noindent\textit{Redator: formato da condição C}
\begin{lstlisting}
Return `answer_text` with the final answer, plus `answer_claim`: the typed claim that the answer asserts. Give it an id of the form W01, set `generated_by` to 'writer', and list in `derived_from` the ids of the analyst claims it came from.
\end{lstlisting}

\noindent\textit{Redator: formato da condição D}
\begin{lstlisting}
Return `answer_text` with the final answer, plus `answer_claim`: the typed claim that the answer asserts. Give it an id of the form W01, set `generated_by` to 'writer', list in `derived_from` the ids of the analyst claims it came from, and carry over the `evidence` snippet.
\end{lstlisting}

\noindent\textit{Condição E (Analista e Redator): como B, acrescida da instrução abaixo}
\begin{lstlisting}
Carry over the `evidence` snippet of the input claim you used.
\end{lstlisting}

\noindent\textit{Condição B' (Analista): como B, acrescida da instrução abaixo}
\begin{lstlisting}
Use `reasoning` to work through the computation before writing the claims.
\end{lstlisting}

\noindent\textit{Condição D' (Analista e Redator): como D, acrescida da instrução abaixo}
\begin{lstlisting}
Before using any input value, check it against its `evidence` snippet; if they disagree, trust the evidence.
\end{lstlisting}

\noindent\textit{Auditor}
\begin{lstlisting}
You are an AUDITOR investigating a pipeline run.

A pipeline (extractor -> analyst -> writer) produced a final answer that may contain a factual error.

Using ONLY the trace below, determine:
- `first_error_stage`: the stage where the first incorrect claim appeared. Use 'none' if you believe the run has no error.
- `root_cause_claim_id`: the id of the claim that started the error chain, or null if the trace has no claim ids.
- `root_cause_claim_text`: a one-line natural-language description of that claim.
- `reported_metric` and `reported_value`: the metric and the incorrect value of that claim.

You have no access to the source document and no ground truth. Judge from internal consistency and from whatever provenance the trace carries.
\end{lstlisting}

Os schemas são modelos Pydantic com \texttt{extra="forbid"}. Os campos de cada claim por condição estão na Tabela~\ref{tab:condicoes}: \texttt{metric} (texto), \texttt{value} (número), \texttt{unit} (texto, com a lista USD\_K, USD\_M, USD\_B, PCT e RATIO só na descrição), \texttt{year} (inteiro), \texttt{claim\_id}, \texttt{source\_id}, \texttt{page}, \texttt{generated\_by}, \texttt{derived\_from} (lista de identificadores) e \texttt{evidence} (texto). O Analista devolve uma lista de claims e, em C e D, o identificador do claim que responde à pergunta. O Redator devolve \texttt{answer\_text} e o claim que a resposta afirma. O relatório do auditor, igual nas quatro condições, tem \texttt{first\_error\_stage} (extractor, analyst, writer ou none), \texttt{root\_cause\_claim\_id}, \texttt{root\_cause\_claim\_text}, \texttt{reported\_metric}, \texttt{reported\_value}, \texttt{confidence} (entre 0 e 1) e \texttt{rationale}.

% !TEX root = ../tcc.tex
% Seção sem número (padrão IEEE). Financiamento vai no \thanks do título, não aqui.
\section*{Agradecimentos}

Aos meus pais, Karina Mirella e Onildo Júnior, minha mais profunda gratidão. Desde cedo, ensinaram-me o valor da educação e me incentivaram a buscar conhecimento e a acreditar em meu potencial. Agradeço por cada esforço, por cada renúncia e por todas as oportunidades que me proporcionaram ao longo da vida. Chegar até aqui é também o resultado de tudo aquilo que fizeram por mim. Por isso, esta conquista não é apenas minha, mas nossa.

À minha noiva, Gabriela Ferreira, minha companheira de vida desde os tempos de colégio, agradeço por todo o amor, apoio e incentivo ao longo dessa caminhada. Tivemos o privilégio de crescer juntos, acompanhar os sonhos um do outro e celebrar cada conquista. Obrigado por acreditar em mim e por sempre me lembrar de que somos capazes de realizar aquilo que sonhamos. Sou imensamente grato por compartilhar a vida com você e por saber que esta conquista é apenas uma entre tantas que ainda construiremos juntos.

Ao meu orientador, Prof. Dr. Cleber Zanchettin, agradeço pela dedicação, pela confiança e pelos ensinamentos compartilhados ao longo do desenvolvimento da pesquisa e deste trabalho.

Aos meus colegas de graduação, com quem compartilhei inúmeros projetos, desafios, aprendizados e conquistas, agradeço por tornarem essa jornada mais leve e enriquecedora. Levo comigo não apenas o conhecimento adquirido, mas também as amizades e as experiências que construímos ao longo desses anos.

Por fim, a todos que, de alguma forma, contribuíram para que eu chegasse até aqui, deixo minha sincera gratidão. Esta conquista carrega um pouco de cada pessoa que fez parte da minha trajetória.

\bibliographystyle{IEEEtran}
\bibliography{referencias}

\end{document}
